\section{Coprimality collapse and the diagonal divisor lift}
\label{sec:divisor-lift}

Throughout this section, \(K:\N^2\to\mathbb C\) has finite support.  All
sums are therefore finite, and no limiting interchange is hidden.

\begin{lemma}[Primitive sign collapse]
\label{lem:primitive-sign-collapse}
For all positive integers \(u,v\),
\begin{equation}
 \boxed{
 \one_{(u,v)=1}\mu(u)\mu(v)=\mu(uv).}
\label{eq:primitive-sign-collapse}
\end{equation}
\end{lemma}

\begin{proof}
If either \(u\) or \(v\) is not squarefree, both sides vanish.  If both are
squarefree and coprime, multiplicativity gives
\(\mu(uv)=\mu(u)\mu(v)\).  If both are squarefree but share a prime, then
\(uv\) is not squarefree, so both sides again vanish.
\end{proof}

The identity is elementary, but its role is structural.  The coprimality
mask and the two signs are not three separate objects; together they are
one product M\"obius coefficient.

\begin{lemma}[Coprimality inversion]
\label{lem:coprimality-inversion}
For all \(u,v\ge1\),
\begin{equation}
 \mu(uv)
 =
 \mu(u)\mu(v)\sum_{d\mid(u,v)}\mu(d).
\label{eq:mu-product-inversion}
\end{equation}
\end{lemma}

\begin{proof}
If \(\mu(u)\mu(v)=0\), both sides vanish.  Otherwise \(u\) and \(v\) are
squarefree, and
\[
 \sum_{d\mid(u,v)}\mu(d)=\one_{(u,v)=1}.
\]
Apply \cref{lem:primitive-sign-collapse}.
\end{proof}

\begin{theorem}[Exact diagonal divisor lift]
\label{thm:diagonal-divisor-lift}
Define
\[
 \cS(K)
 =
 \sum_{\substack{u,v\ge1\\(u,v)=1}}
 \mu(u)\mu(v)K(u,v).
\]
Then
\begin{align}
 \cS(K)
 &=
 \sum_{u,v\ge1}\mu(uv)K(u,v)
\label{eq:plane-as-product-mu}\\
 &=
 \sum_{d\ge1}\mu(d)
 \sum_{\substack{x,y\ge1\\(x,d)=(y,d)=1}}
 \mu(x)\mu(y)K(dx,dy).
\label{eq:exact-diagonal-lift}
\end{align}
\end{theorem}

\begin{proof}
The first equality is \cref{lem:primitive-sign-collapse}.  Insert
\eqref{eq:mu-product-inversion} and interchange finite sums:
\[
 \cS(K)
 =
 \sum_{d\ge1}\mu(d)
 \sum_{x,y\ge1}\mu(dx)\mu(dy)K(dx,dy).
\]
Only squarefree \(d\) contribute.  For such \(d\),
\[
 \mu(dx)
 =
 \mu(d)\mu(x)\one_{(x,d)=1},
\]
with the same identity for \(y\).  The two factors \(\mu(d)\) square to
one, leaving \eqref{eq:exact-diagonal-lift}.
\end{proof}

\begin{remark}[The diagonal variable is new]
The divisor \(d\) in \eqref{eq:exact-diagonal-lift} is introduced by
coprimality inversion.  It is not the row cofactor \(d_m\), the anchor
content \(t\), or the erasure content \(q\) of
\cref{sec:primitive-interface}.  Conflating these variables would change
both the support and the sign.
\end{remark}

\begin{remark}[Why the lift helps]
For a fixed \(d\), the two variable signs in
\eqref{eq:exact-diagonal-lift} separate.  Their only modification is the
fixed local condition \((x,d)=(y,d)=1\).  The price is a sum over diagonal
dilations \(K(dx,dy)\).  Low \(d\) will be handled by the prime number
theorem; large \(d\) will be an explicit tail.
\end{remark}

\subsection{Removing one primitive direction}

The physical erasure plane excludes the anchor direction \((b,a)\), where
\((a,b)=1\).  For any kernel \(K\), define
\[
 K^\circ=K-K(b,a)e_{b,a},
 \qquad
 e_{b,a}(u,v)=\one_{(u,v)=(b,a)}.
\]
Then
\begin{equation}
 \cS(K^\circ)
 =
 \cS(K)-\mu(a)\mu(b)K(b,a).
\label{eq:anchor-point-removal}
\end{equation}
Thus the unequal-residual condition is exactly one primitive-plane site
correction.  Its value may aggregate several underlying labels and is not
called one physical atom.  For the literal TPC--70 erasure kernel the site
is already zero; the formula is useful only if one first works with an
enlarged bookkeeping kernel.

\begin{lemma}[Diagonal visibility of the removed point]
\label{lem:anchor-diagonal-visibility}
Let \(K^{[d]}(x,y)=K(dx,dy)\).  The removed point \((b,a)\) appears in
\((K^\circ)^{[d]}\) only for \(d=1\).  For every \(d>1\),
\[
 (K^\circ)^{[d]}=K^{[d]},
\]
and hence
\[
 \cV_\square\bigl((K^\circ)^{[d]}\bigr)
 =
 \cV_\square(K^{[d]})
\]
For \(d=1\),
\[
 \cV_\square(K^\circ)
 \le
 \cV_\square(K)+4|K(b,a)|.
\]
\end{lemma}

\begin{proof}
The point survives the \(d\)-dilation only if \(d\mid a\) and \(d\mid b\).
Since \((a,b)=1\), this forces \(d=1\).  A point mass has four nonzero
mixed differences, counted with multiplicity at the axes when necessary,
and hence mixed variation at most four times its amplitude.
\end{proof}
